\section{The molecular layer, the trace metals and element transport}
\label{sec:mol}

Everything in this section is \emph{off} by default. A run acquires the
molecular layer by asking for it (\texttt{Molecular chemistry: True}, which
also switches \texttt{Molecular base} on so that the base particle count and
the chemistry describe one gas), the metals by the presence of \texttt{metals.inp} or of elemental
abundances in a handoff file, and the element transport by
\texttt{He\_diffusion}. An atomic, metal-free run never enters any of the
code described here.

\subsection{\texorpdfstring{The H$_2$/He reaction network}{The H2/He reaction network}}
\label{sec:mol:network}

\texttt{mol\_rates.f90} carries the reaction rate coefficients of
\citet{Koskinen2022} Table 1 as published, in cgs
(cm$^{3}$\,s$^{-1}$, and cm$^{6}$\,s$^{-1}$ for the two three-body entries
R13 and R15; R13 is multiplied by the total heavy-particle density and R15 by
a collider-resolved one, below).
EXHALE carries one temperature, so the heavy-particle and the electron
temperature of that table are the same variable. The photorates P1--P5 of
the paper are marked ``SC'' there, computed from cross sections, stellar
flux and columns; in EXHALE they are the integrals of
Sec.~\ref{sec:radiation:rates} and are not part of this module.

One entry is deliberately \emph{not} transcribed. R12, the thermal
dissociation of H$_2$, is built by detailed balance from the R15
recombination,
\begin{equation}
k_{12}(T) = \frac{k_{3b}(T)}{K_{\rm eq}(T)},
\label{eq:mol:r12}
\end{equation}
with $K_{\rm eq}$ the ${\rm H}+{\rm H}\rightleftharpoons{\rm H_2}$
equilibrium constant built from the same internal partition function
$Q_{\rm H_2}(T)$ the caloric equation of state uses
(Sec.~\ref{sec:mol:caloric}), because the two Table-1 entries are the same
channel in opposite directions and their source measurements do not overlap
in temperature. The third body of R13 is the total gas-particle density
excluding electrons, its source (Miller et al.\ 1968) being a single
coefficient with no collider dependence measured. R12 and R15 are
collider-resolved (corrected 2026-09-17, item L7g): the association rate is
the sum $\sum_M k_1(M)\,n_M$ over $M = {\rm H_2}$, H and He with the
published coefficients of \citet{CohenWestberg1983}, carried in the code as
an H$_2$-equivalent third-body density so that R12 keeps following R15 by
\eqref{eq:mol:r12}.  Applying the $M={\rm H_2}$ coefficient to the total
density, which is what the code did before, over-counts the association: at
the base of the LHS 1140 b hot-Uranus column it made the R15 rate 22 per cent
too large and the molecular chemical heat 21.6 per cent too large.  No
measurement of He as a third body for this reaction exists, so He is given
argon's coefficient, and that entry remains an upper bound: a monatomic third
body has no internal modes to take up the released energy. Two
further limits in a helium-dominated gas, the R16 and R18 differences
of order unity to ten with \citet{GarciaMunoz2025}, are recorded at the code
site rather than patched.

R19, ${\rm HeH^+}+{\rm H}\to{\rm H_2^+}+{\rm He}$, is not the Table~1
constant $9.1\times10^{-10}$\,cm$^3$\,s$^{-1}$, a measurement near 300~K that
held at every temperature of the wind is an extrapolation. The code carries
the thermal rate of \citet{Esposito2015}, quantum close-coupling cross
sections for the lowest states and quasi-classical trajectories for the rest,
Boltzmann-averaged over all 178 rovibrational states of HeH$^+$, in the
authors' fit $k=\sum_{i=0}^{6}a_i(\log_{10}T)^i$ (maximum error 2.5\%;
read in the accepted manuscript). It is $1.09\times10^{-9}$ at 300~K,
$1.52\times10^{-9}$ at 1000~K and $2.00\times10^{-9}$ near 6000~K, close to
the Langevin value $2.09\times10^{-9}$, and lies within 1--6\% of the exact
quantum reactive rate of \citet{DeFazio2014} from 100 to 2000~K; above 1000~K
it is 1.7--2.2 times the Table~1 constant. The fit is held at its 10 and
15000~K ends outside that range, the authors call the rate questionable above
about $10^4$~K, and the thermal average assumes HeH$^+$ internal states in
equilibrium at the gas temperature. On the hot-Uranus fixtures the change
lowers $n({\rm HeH^+})$ by up to 46\% and leaves the gas temperature,
velocity and density within $4\times10^{-5}$ (Update log stage 3,
section 64).

\textbf{Two entries of Table~1 are not transcribed as printed.} The first is
R17, ${\rm He^+}+{\rm H_2}\to{\rm H^+}+{\rm H}+{\rm He}$, the dissociative
charge transfer. Table~1 gives it as $10^{-9}\exp(-5700/T)$
\citep{Moses2000}, an Arrhenius form with a $0.49$~eV barrier, and that form
alone is $2\times10^{4}$ below the directly measured rate at 300~K and 200
below it at 400~K. The reaction has two mechanisms: at thermal energy it
proceeds by tunnelling out of a long-lived ${\rm He^+}$--${\rm H_2}$ complex,
because reactants and products do not correlate adiabatically, and only above
$\sim400$~K does an over-barrier branch take over
\citep{Boehringer1986}. The drift-tube measurements give the two-body total
as $k_2=1.1\times10^{-13}(300/T)^{0.24}$~cm$^3$\,s$^{-1}$ over 18--408~K
\citep{Boehringer1986}, consistent with $(1.1\pm0.1)\times10^{-13}$ at 330~K
and $(1.5\pm0.15)\times10^{-13}$ at 78~K \citep{Johnsen1980}. The code
therefore carries the sum of the two mechanisms, with the tunnelling branch
taken as the measured total less the radiative branch the network carries
separately as R23, so that the coded channels sum to the measured total:
\[
   k_{\rm R17}(T)=\Bigl[1.1\times10^{-13}\,(300/T)^{0.24}-k_{\rm R23}\Bigr]
                  +10^{-9}\exp(-5700/T).
\]
Against the only measurement of the rising branch \citep[][their Fig.~3,
carried to an effective temperature of 700~K]{Johnsen1980} the sum is inside
a factor two from 300 to 700~K. The three-body channel both papers measure,
$k_3=1.6\times10^{-30}(100/T)^{1.27}$~cm$^6$\,s$^{-1}$, is omitted because
$k_3n=6\times10^{-19}$~cm$^3$\,s$^{-1}$ at 1000~K and $n=7\times10^{12}$~cm$^{-3}$,
five decades below the two-body rate.

\textbf{The second is R20.}
Table~1 gives ${\rm He^+}+{\rm H_2}\to{\rm HeH^+}+{\rm H}$ as
$4.2\times10^{-13}$~cm$^3$\,s$^{-1}$, citing \citet{Schauer1989}. That paper
measures the radiative
(${\rm He^+}+{\rm H_2}\to{\rm He}+{\rm H_2^+}+h\nu$, Table~1's R23) and the
dissociative (${\rm He^+}+{\rm H_2}\to{\rm He}+{\rm H}+{\rm H^+}$, R17)
charge transfer over $15\le T\le40$~K and no HeH$^+$ channel; it states that
the HeH$^+$ channel ``apparently does not become allowed until the collision
energy approaches 9~eV'', and its H$_3^+$ signal is the \emph{sum} of the
radiative and the HeH$^+$ channels, which bounds the latter at
$\le1.0\times10^{-14}$ -- 42 times below the value Table~1 assigns it. R20 is
therefore not carried. The HeH$^+$ source the network uses instead is the one
Table~1 omits, ${\rm H_2^+}+{\rm He}\to{\rm HeH^+}+{\rm H}$ at
$3.0\times10^{-10}\exp(-6717/T)$ \citep{Black1978}, whose coefficient is a
Maxwellian average over a thermal population of H$_2^+$ vibrational levels.

Measured on LHS~1140~b at ${\rm He/H}=2.13$ with the H$_2$ row written out
reaction by reaction, the repair rewires the base layer's HeH$^+$ and H$^+$
budgets -- at the first cell $n({\rm HeH^+})$ falls by a factor $487$ at
1023~K and $8\times10^4$ at 629~K, and $n({\rm He^+})$ rises by $1.11$ and
$4.4$ -- and leaves the net H$_2$ destruction there within a per cent,
because the helium ions of that layer are sink-limited. The same measurement
shows what sets the H$_2$ of a helium-rich base: the helium-ion and proton
channels, with the thermal dissociation R12 nine decades below the first
cell's total H$_2$ loss at 1023~K and twenty-two decades below it at 629~K,
and nowhere in that layer reaching 0.4 per cent of the loss. The helium-ion channels, not the temperature and not the photons, set
the H$_2$ content of a helium-rich base in this network.

The eight balance rows the network makes are
\eqref{eq:ion:h2row}--\eqref{eq:ion:hehprow} together with the H$^+$,
He$^+$, He$^{2+}$ and He\,$2^3$S rows of Sec.~\ref{sec:ionization:cellsystem}.

\subsection{\texorpdfstring{H$_2$ photoabsorption resolved into channels}{H2 photoabsorption resolved into channels}}
\label{sec:mol:h2photo}

\texttt{h2\_photo\_channels.f90} splits the one measured cross section
$\sigma_{\rm H_2}(E)$ into four mutually exclusive final states whose sum is
identically $\sigma_{\rm H_2}$, so no consumer of the opacity changes and
the total H$_2$ destruction rate is what it was; what changes is which
fragments the source terms are told to make.
\begin{align}
\sigma_N &= f_n(E)\,\sigma_{\rm H_2}(E)
  &&{\rm H_2}+h\nu\to{\rm H}+{\rm H}\ \ (33\text{--}41\,{\rm eV\ window}),
  \label{eq:mol:chanN}\\
\sigma_M+\sigma_S+\sigma_D &= \left[1-f_n(E)\right]\sigma_{\rm H_2}(E)
  &&\text{the three ionizing channels},
  \label{eq:mol:chanMSD}
\end{align}
with $f_n$ the neutral fraction of \citet{Chung1993} Table 1, carried as
a \emph{fraction} so that it rides on this code's $\sigma_{\rm H_2}$
normalization rather than importing theirs. Inside the ionizing part the
split follows the \citeauthor{Chung1993} Table II ratio column
$r(E) = \sigma({\rm H^+})/\sigma({\rm H_2^+})$, a ratio of time-of-flight
detector signals and therefore a ratio of \emph{events}, one count per event
and not one per proton;
\texttt{frac\_H2\_dissociative\_ionization} is that column converted to
$\sigma({\rm H^+})/[\sigma({\rm H^+})+\sigma({\rm H_2^+})]$ and tabulated at
69 energies from 18.076 to 124\,eV, linearly interpolated, zero at and below
threshold and held flat above 124\,eV. The thresholds are 15.4\,eV (M,
leaves H$_2^+$), 18.076\,eV (S, leaves H\,$+$\,H$^+$), 51.4\,eV (D, two
protons), and none of its own for N. The D channel is resolved by default
(\texttt{H2 double ionization: chung80}, the 20 per cent of the H$^+$
events at and above 80\,eV that \citet{Chung1993} state, ramped linearly from
the 51.4\,eV threshold, the ramp not being in the source);
\texttt{yan\_rho} builds it from the double-to-single ratio of
\citet{Yan1998} instead (at 110\,eV \texttt{chung80} gives 15 per cent more
double ionization than \texttt{yan\_rho}), and
\texttt{off} folds it back into S, which reproduces the code before the
channel existed. The N window is resolved by default as well
(\texttt{H2 neutral dissociation: True}).

The energy of one absorbed photon is accounted channel by channel by
\texttt{h2\_channel\_energy\_recipients}, the single statement of what each
channel charges, and the recipients of each sum to $h\nu$ once. Two entries
are fixed with no free parameter: the double channel's vertical threshold
51.4\,eV \citep{Yan1998} exceeds the asymptotic product
energy $D_0+2I({\rm H}) = 31.675$\,eV by exactly the kinetic energy release
of the two receding protons, 19.725\,eV; and the neutral window leaves both
fragments in $n=2$ (the Q2 $^1\Pi_u$ state whose fluorescence curve overlaps
this cross section), so $2E(n{=}1\to2)$ departs as prompt Ly$\alpha$ and
two-photon continuum and is not heat.

\subsection{Lyman-Werner photodissociation}
\label{sec:mol:lw}

\texttt{lyman\_werner.f90} adds the Solomon process,
${\rm H_2}(X)+h\nu(11.2\text{--}13.6\,{\rm eV})\to{\rm H_2}(B,C)\to{\rm H}+{\rm H}$,
which the \citeauthor{Koskinen2022} Table-1 network leaves out (their photorates start
at the 15.4\,eV edge). The band is 912--1201\,\AA{}: 912\,\AA{} is the H
Lyman edge, and 1201\,\AA{} is both the blue edge of band B2 and the red end
of the Lyman and Werner line list the shielding table is built from. \citeauthor{Draine1996}'s own 912--1110\,\AA{} is not used, because a planetary base runs
at 700--3200\,K where the vibrationally excited levels pump in lines longward
of 1110\,\AA{} that carry about a third of the dissociations at saturation.
One band, one line list, one normalization.

The optically thin rate is formed from the band-integrated \emph{energy}
flux at the planet and the mean photon energy of a flat $F_\lambda$ band,
$\langle h\nu\rangle = 2hc/(912+1201\,\mbox{\AA}) = 11.7354$\,eV:
\begin{equation}
k_{\rm LW}^{\rm thin} = \sigma_{\rm LW}\,
   \frac{\xi_{\rm day}F_{\rm LW}}{\langle h\nu\rangle}.
\label{eq:mol:klw}
\end{equation}
The band-mean $\sigma_{\rm LW}$ is not the effective cross section
$2.263\times10^{-18}$\,cm$^{2}$ that \citet{Draine1996} Table 1 and 2 imply
for a flat-$F_\lambda$ field restated per 912--1201\,\AA{} photon: the rate
is the level-resolved table of \texttt{h2\_self\_shielding\_table.f90},
normalized per photon of the same band, and that number is only the limit it
approaches in optically thin interstellar conditions.

That table is a line-by-line calculation with \emph{overlapping} lines, and
it stores three different counts because three are needed:
$\sigma_{\rm diss}$, dissociations per unit incident band photon fluence,
which carries the self-shielding, the line overlap and the trapping of the
fluorescent photons and sets the rate and the fragment heating;
$p_{\rm single} = D/(D+\sum A_{ul})$, dissociations per pump with every
fluorescent photon free to escape, which sets the vibrational heat return
because a trapped decay plus its re-absorption moves one molecule down and
another up and deposits nothing; and
$p_{\rm eff} = D/(D+\sum A_{ul}P_{\rm loss})$, dissociations per photon
removed from the beam, which is the band photon ledger. They are related by
$\sigma_{\rm diss} = p_{\rm eff}\,\sigma_{\rm pump}$, and the beam's own loss
is served from the same table, so the photons the beam loses and the
dissociations the rate spends are one absorption with one normalization.
Line overlap is not a correction at these columns: at
$n_{\rm H}=10^{13}$\,cm$^{-3}$ the surviving pumping relative to a
one-line-at-a-time calculation is 0.45 at $N_{\rm H_2}=10^{20}$\,cm$^{-2}$
and 900\,K, and $4.8\times10^{-5}$ at $4.4\times10^{21}$\,cm$^{-2}$ and
2700\,K, which is where the base of a molecular layer sits. The rate enters
\eqref{eq:ion:h2row} as $g_{\rm LW}$ and is zero when no band flux is
supplied.

\subsection{The caloric equation of state}
\label{sec:mol:caloric}

The caloric equation of state that gives the bound H$_2$ its rovibrational
heat capacity is \eqref{eq:caloric} of Sec.~\ref{sec:hydro:eos}, where it is
stated with its validity and its \texttt{Caloric EOS: monatomic} switch.
\texttt{caloric\_eos.f90} supplies the four maps the hydrodynamics needs,
$(\rho,\rho e)\to T$, $(\rho,\rho e)\to p$, $(\rho,p)\to\rho e$ and
$(\rho,p)\to\gamma_{\rm eff}$, and $\gamma_{\rm eff}$ itself for the wave
speeds.

\subsection{Carrier transport}
\label{sec:mol:carrier}

Everywhere else in the code the composition of a cell is the root of a local
algebraic system. A \emph{carrier} is a species the code does not fix that
way: one whose chemical time can exceed the flow time, so that its density
in a cell is what the flow brought there and not what the cell's own
equilibrium would give. \texttt{diffusive\_photochemistry.f90} solves, for
the carriers, the continuity equation of the neutral carriers H$_2$, OH,
H$_2$O and CO (and, with \texttt{Ionization transport}, the proton),
\begin{equation}
\frac{\partial n_i}{\partial t}
 + \frac{1}{r^{2}}\frac{\partial}{\partial r}
   \left[r^{2}\left(n_i v + J_i\right)\right] = P_i - L_i,
\qquad
J_i = -n_{\rm tot}\left(D_i+K_{zz}\right)\frac{\partial f_i}{\partial r}
         - n_i D_i\left(\frac{1}{H_i}-\frac{1}{H_{\rm atm}}\right),
\label{eq:mol:carriereq}
\end{equation}
with $f_i = n_i/n_{\rm tot}$. Between them these four carry every H nucleus
that is not atomic or in a molecular ion and every O and C nucleus that is
not a free atom. H$_2^+$, H$_3^+$, HeH$^+$ and O($^1$D) are left local; on
the HD\,209458\,b molecular example the three molecular ions together carry
at most $1.2\times10^{-9}$ of the H nuclei, and the H-nucleus budget is
closed against the element total rather than against a sum of transported
species, so whatever they hold is simply not available to the carriers. All
four transported species are neutral, so the ambipolar field term of the
general drift-diffusion equation is absent rather than neglected.

$D_i$ is Blanc's law over the background carriers (H\,{\sc i}, H\,{\sc ii},
H$_2$, H$_2^+$, H$_3^+$, He\,{\sc i}, He\,{\sc ii}, He\,{\sc iii},
HeH$^+$), each pair from the rigid-sphere coefficient of \citet{Banks1973}. The ion-neutral polarization enhancement is not added: the
transported species only exist where the gas is neutral (electron fraction
per H nucleus at most $10^{-7}$ over the cells that hold H$_2$), so the ion
pairs carry that share of the Blanc friction sum, and for H$_2$O the
induced-dipole formula would be the wrong one anyway, a permanent dipole of
1.85\,D dominating the capture rate. $K_{zz}$ is read from
\texttt{kzz\_cell} and nowhere else; that array is filled once, from the
scalar \texttt{He\_Kzz} or from the lower-atmosphere profile's own $K_{zz}$
column. Settling is $G_i = (m_i-\bar m)g/(kT)$ with $\bar m = \rho/n_{\rm tot}$
and $g$ from the same potential the hydrodynamics uses; thermal diffusion is
not carried ($\alpha_T=0$), there being no evaluated $\alpha_T$ for H$_2$O,
OH or CO in an H$_2$/He bath.

$P_i-L_i$ is not written in this module: it is obtained by calling the very
rows the local equilibrium solve uses, \texttt{mol\_heh\_rows} row 4 for
H$_2$ and \texttt{oxygen\_carrier\_rows} for OH and H$_2$O, at the trial
densities with the background frozen, and the Jacobian is a forward
difference of those same calls, so a rate that changes there changes here.
CO has no chemical source and only its transport moves it, with a ceiling
from \texttt{limit\_to\_element\_budget} that keeps it from carrying the
oxygen and the carbon out of the wind, where the molecule cannot survive and
where the O\,{\sc i} and C\,{\sc ii} lines are the coolant.

The discretization is cell-centered $f_i$, face-centered fluxes on
\texttt{r\_edg}, one backward-Euler step per call, Newton with a
backtracking line search, block tridiagonal with $4\times4$ blocks. Faces 0
and $N$ carry zero diffusive flux; the drift term switches from central to
donor cell on a P\'eclet test. The element reservoirs at the base are
Dirichlet, imposed through \texttt{melem\_ab}, the base density and the
equation of state, but the \emph{partition} of each element among its
carriers is not imposed anywhere, which is what makes the comparison against
a photochemical handoff a test rather than a tautology.

The material advection of a carrier is \emph{not} in these rows. It is the
divergence of the hydrodynamic face mass flux, $F_\rho(j)Y_c(j)/m_c$, taken
with the mass row itself inside the Runge-Kutta stages \eqref{eq:ssprk3} on
the same faces, areas, volumes and time step
(\texttt{species\_face\_flux.f90}), so a cell
can only lose the carrier the mass row says it loses and a uniform partition
is preserved to round-off. That this half must ride on the \emph{face} mass flux is a numerical
statement about the front: a first-order donor-cell difference on the cell
velocity carries a numerical diffusivity $|v|\Delta r/2$, which on the
500-cell hot Uranus grid is 8.5, 20 and 22 times the molecular $D$ of H$_2$
at 1.15, 1.30 and $1.50\,R_{\rm p}$. The key is
\texttt{Molecular carrier transport}, default off.

\subsection{The oxygen option}
\label{sec:mol:oxygen}

With \texttt{thereis\_oxychem} (default off) three neutral carriers OH,
H$_2$O and CO join the network. \texttt{oxygen\_rates.f90} carries the rate
coefficients, each transcribed from a named publication and each checked
against two independent ports of the same literature that are on this
machine (the Photochem \texttt{zahnle\_earth.yaml} network and the VULCAN
\texttt{NCHO\_photo\_network.txt}), every disagreement being recorded at the
coefficient. \texttt{water\_photolysis.f90} supplies the photolysis,
\[
{\rm H_2O}+h\nu\to{\rm OH}+{\rm H},\quad
\to{\rm H_2}+{\rm O}(^1{\rm D}),\quad
\to{\rm O}+{\rm H}+{\rm H};
\qquad
{\rm OH}+h\nu\to{\rm O}+{\rm H},
\]
in the four far-ultraviolet bands of Sec.~\ref{sec:radiation:spectrum}
(the balance rows these rates enter are \eqref{eq:ion:ohrow} and
\eqref{eq:ion:h2orow}), whose
edges are fixed by the branching ratios rather than chosen: the three-body
branch opens only in the interval containing Ly$\alpha$ and is 12 percent
there, so Ly$\alpha$ needs its own band on branching grounds alone. The first
branch alone carries 54.5 percent of the H$_2$O loss at the HD\,189733\,b
base, which makes the OH/H$_2$O cycle catalytic rather than a one-way sink.
\texttt{co\_photodissociation.f90} destroys CO on the same Lyman-Werner beam,
using a band-mean dissociation cross section times the band photon flux and
the \citet{Visser2009} shielding function
$\Theta(N_{\rm CO},N_{\rm H_2})$ rather than their $k_0\Theta$ form, because
EXHALE has a band flux and not an interstellar field and $k_0$ is a statement
about a particular field. The other CO destruction channel is the He$^+$
charge transfer
${\rm He^+}+{\rm CO}\to{\rm C^+}+{\rm O}+{\rm He}$ at the Langevin value
$1.6\times10^{-9}$\,cm$^{3}$\,s$^{-1}$, which appears in the He$^+$ row of
\eqref{eq:ion:h2prow}'s system and whose $+2.2117$\,eV appears in heating
channel 18, one event counted once.

\subsection{The lower-atmosphere handoff}
\label{sec:mol:handoff}

The base of the escape model sits at a stated pressure level, and three
routes supply the state there.

\paragraph{The analytic column.} \texttt{lower\_column.f90} integrates the
hypsometric relation from the 1-bar level up to the base pressure (default
1\,$\mu$bar), isothermal at $T_{\rm eq}$ with spherical gravity, using the
mean molecular weight of the chemical-equilibrium H$_2$/H/He partition of
\citet{Visscher2006} as fitted by \citet{Koskinen2022} eqs.~(11)--(13),
\begin{equation}
q_{\rm H_2} = \frac{1.9845+10^{u}-\sqrt{10^{u}\left(3.9690+10^{u}\right)}}{2.3670},
\qquad
u = -\frac{23672}{T}-\log_{10}p\,[{\rm bar}]+6.2645,
\label{eq:mol:qh2}
\end{equation}
with $q_{\rm He}$ from the elemental He/H and $q_{\rm H} = 1-q_{\rm H_2}-q_{\rm He}$.
It returns the radius of the base level and the base mixing ratios, and
also a fully atomic bracket, because in chemical equilibrium
$q_{\rm H_2}(1\,\mu{\rm bar})$ stays large up to $T\sim2000$\,K while
photochemistry dissociates H$_2$ there: a large $q_{\rm H_2}$ on a hot planet
is the signal that a photochemical handoff is needed. The validation case is
the \citet{Koskinen2022} Model A Uranus-mass planet at 0.05\,au, for which
the column gives $r_0(1\,\mu{\rm bar}) = 1.34\,R_{1\rm bar}$ with
$q_{\rm H_2}\simeq0.84$.

\paragraph{\texttt{base.inp}, the scalar handoff.} An optional file of
``key value'' lines, written by \texttt{src/utils/run\_lower.py} or
\texttt{src/utils/vulcan\_to\_base.py}. Every key belongs to one of five
categories and the category states what the value may do to the wind:
\emph{provenance} (comments only); \emph{EOS boundary}
(\texttt{T\_base}$\to T_0$, \texttt{r\_base}$\to R_0$, \texttt{p\_base}, and
\texttt{q\_H2\_base}, the H$_2$ volume mixing ratio of the base gas);
\emph{elemental reservoir} (\texttt{HeH\_base} and
\texttt{<El>\_H\_base} for the ten elements, overriding \texttt{metals.inp}
for that element); \emph{initial guess} (no key today); and \emph{boundary
constraint} (\texttt{Kzz\_base}$\to$\texttt{he\_kzz}). \texttt{q\_H2\_base}
is the $q_{\mathrm{H_2}}$ of \eqref{eq:ntotbc} and is imposed on the H$_2$
partition of the inflowing ghost species as well; it is refused above
$0.5/(0.5+{\rm He/H})$. Species mixing ratios other than \texttt{q\_H2\_base}
stay comments: no part of the code consumes them, so presenting them as
physics would be false. A \texttt{base.inp} that states \texttt{p\_base} fixes the base level and
$n_0$ is derived from it as in Sec.~\ref{sec:intro:units}, instead of being
read from \texttt{input.inp}.

\paragraph{\texttt{Lower atmosphere profile}, the tabulated handoff.}
\texttt{lower\_atmosphere\_profile.f90} reads the lower model's solution
over an \emph{interval} of pressure that reaches above the level where
EXHALE places its base. It supplies the state at the matching level
$p_{\rm match}$ ($T$, $r$, $q_{\rm H_2}$, and the elemental ratios He/H and
El/H), which is what the base equation of state and the elemental reservoirs
are built from, and $K_{zz}$ as a \emph{profile}, because the homopause,
the only consequential thing $K_{zz}$ does, is where $K_{zz}=D_{12}$, a
property of two profiles, and the molecular coefficient rises by nearly five
decades between the base and $1.2\,R_{\rm p}$. Every intensive quantity is
interpolated linearly in $\log p$, the variable both models solve on;
extrapolation is never performed. The values go in through exactly the doors
the scalar keys use, so nothing in the base equation of state changes; $n$
and $\rho$ carried by the file are \emph{not} imposed, the base density
remaining $n_0$.

\subsection{Trace metals}
\label{sec:mol:metals}

\texttt{metals.inp} (\texttt{metals\_input\_read.f90}) sets the elemental
abundances $n_X/n_{\rm H}$ by number, one ion per line with the neutral label
naming the element; a second statement of the same element stops the run,
since two abundances for one element are two answers to one question. The
same file also carries \texttt{cno\_cool}, \texttt{eos\_metals},
\texttt{cx\_full} and \texttt{cx\_O2p\_H}. The ionization of each element is
solved in the same coupled system as hydrogen and helium
(\texttt{System\_HeH\_metals.f90} and its triplet and molecular
counterparts), through the rows \eqref{eq:ion:metalrows}; the code carries no
parallel decoupled metal solver.

\texttt{eos\_metals} (default \emph{on}) puts the metals into the bulk gas
budget consistently: their mass in $\rho$ and hence in $\mu$, $v_0$, $p_0$;
their electrons in $n_e$; and their nuclei in $n_{\rm tot}$ and in the ghost
pressure, matching the charge balance the ionization systems already use.
The legacy trace approximation (metals in the ionization and cooling but not
in the bulk gas) is \texttt{eos\_metals 0}. A metals-off run is identical
either way, every metal sum vanishing.

\subsection{Binary element transport}
\label{sec:mol:diffusion}

\texttt{binary\_element\_diffusion.f90} treats the gas as two components
moving through each other with the single mass-averaged velocity of the
hydrodynamics: component 1 is the hydrogen carriers together with the trace
metals and the heavy nuclei bound in the molecular carriers, and component
He is helium in all stages. Since $\rho_1+\rho_{\rm He}=\rho$ exactly, the
two diffusive mass fluxes close, $J_1=-J_{\rm He}$, and the composition
update creates no mass. The transported variable is the helium \emph{mass}
fraction $X = \rho_{\rm He}/\rho\in[0,1]$, which is bounded at both ends
unlike the ratio $n_{\rm He}/n_{\rm H}$ that the earlier trace kernel
carried and that diverges exactly in the helium-rich limit this operator
exists for. The equation is
\begin{equation}
\frac{\partial(\rho X)}{\partial t}
 + \frac{1}{r^{2}}\frac{\partial}{\partial r}
   \left[r^{2}\left(\rho X v + J\right)\right] = 0,
\qquad
J = -\rho\left(D_{12}+K_{zz}\right)\frac{\partial X}{\partial r}
    - \rho D_{12}\,G\,X(1-X),
\label{eq:mol:hediff}
\end{equation}
\begin{equation}
G = \frac{\left(m_{\rm He}-m_{c1}\right)m_{\rm H}g
          -\left(\bar Z_{\rm He}-\bar Z_{c1}\right)eE}{kT}
   + \alpha_T\frac{{\rm d}\ln T}{{\rm d}r}
   - \frac{{\rm d}\ln\psi}{{\rm d}r}\qquad[{\rm cm^{-1}}].
\label{eq:mol:hedriftG}
\end{equation}
The masses and charges in $G$ are \emph{carrier} quantities:
$\psi = n_{1c}/n_{\rm H}$ is the number of collision partners each hydrogen
nucleus is spread over (1 atomic, $1/2$ fully molecular) and
$m_{c1}=m_1/\psi$. The gradient coefficient collapses to $\rho D_{12}$ for
any $\psi$, so the Fickian term is the same in the molecular and the atomic
region; the piece of the mole-fraction gradient the chemistry carries does
not cancel and rides in $G$ as $-{\rm d}\ln\psi/{\rm d}r$, which vanishes
identically wherever $\psi$ is constant. The ambipolar field is computed and
not assumed, following \citet{Koskinen2013} section 2.1,
\begin{equation}
eE = -\frac{1}{n_e}\frac{{\rm d}(n_ekT)}{{\rm d}r}
   = -kT\frac{{\rm d}\ln(n_eT)}{{\rm d}r},
\label{eq:mol:ambipolar}
\end{equation}
in the logarithmic form, which is exact for the barometric stratification of
the base and the same second-order central difference elsewhere. It is
gated by \texttt{he\_ambipolar} (default on); $\alpha_T$ is 0 by default
(\texttt{He\_alphaT}).

The equation is split. The advection of the element mass fraction is a
divergence of the same face mass flux the density rides on and is carried
inside the Runge-Kutta stages \eqref{eq:ssprk3}, so each element's nucleus total over the
domain changes only by its boundary fluxes to round-off. What is left for
the implicit step here is
$\rho(X^{\rm new}-X^{\rm adv})/\Delta t = -r^{-2}\partial_r(r^{2}J)$, taken
at the composition the stages advected; the operator advects only when it is
\emph{given} the face mass flux, which happens in the fixed-wind relaxation
that has no stages to ride on.

The discrete operator is one backward-Euler step with frozen coefficients,
faces carrying $r^{2}$ areas, Newton on a tridiagonal Jacobian (the step is
nonlinear in $X$ through the product $X(1-X)$). The gradient term is central;
the drift term keeps a P\'eclet hybrid, central where
$|B|\Delta r \le 2(A+E)$ and donor cell otherwise. The two factors of
$X(1-X)$ come from \emph{opposite} sides of the face, because the drift is a
counter-flow and a donor-cell rule must take each element's mass fraction
from the cell that element leaves:
\begin{equation}
J_{\rm drift}(f) =
\begin{cases}
 -B_f\,X(j{+}1)\left[1-X(j)\right], & B_f\ge0\ \ (\text{helium inward}),\\
 -B_f\,X(j)\left[1-X(j{+}1)\right], & B_f<0.
\end{cases}
\label{eq:mol:driftface}
\end{equation}
A donor at $X=0$ then sends nothing and an acceptor at $X=1$ receives
nothing, so a cell on either end of the axis cannot be pushed past it; with
the P\'eclet switch the same holds in the central branch, and the row
argument (time term, gradient flux, upwind advection and drift all of one
sign at a cell that first touches an end) gives $0\le X\le1$ on the residual
itself rather than through a comparison principle. Measured, the worst excursion outside $[0,1]$ before the assertion clip is
$-1.05\times10^{-1}$.

\texttt{project\_elements} then rebuilds the species vector on the new
element totals: helium-bearing species are scaled by
$r_{\rm He}=n_{\rm He}^{\rm new}/n_{\rm He}^{\rm old}$, component-1 species
by $r_1$, and a species carrying a nucleus of each (HeH$^+$) by
$\min(r_1,r_{\rm He})$, the nuclei this undercounts for the element with the
larger factor being deposited into that element's neutral ground species.
Every operation is a multiplication by a nonnegative factor or an addition,
so no negative intermediate can arise and both element totals are met
exactly. The invariant is that the species handed back carry the mass the
species handed in carried, to round-off, which requires $r_1$ to be the ratio
of the two component-1 \emph{masses} read from the cell. The result is the
initial guess for \texttt{ioniz\_eq}, which owns the split \emph{within} an
element while this step owns the element totals.

\paragraph{The base of the operator, and what crosses it.}
The gradient and drift coefficients are built on the faces 1 to $N-1$, and
faces 0 and $N$ are left at zero, so \textbf{the diffusive flux at the base
face is zero by construction and not by the state}: an element crosses either
end of the column with the gas alone.  The flux the Dirichlet reservoir
composition drives sits at the first solved face, which is face 1 while the
operator solves rows 2 to $N$ against the reservoir, and face 0 where a
closed base makes cell 1 an unknown.  The record
\texttt{element\_base\_flux} carries both faces rather than leaving either
to be assumed: the advective and the diffusive half at the base face and at
the first solved face, in ${\rm g\,cm^{-2}\,s^{-1}}$ positive outward
(toward increasing $r$), together with the index of the first solved face and
a count of the operator evaluations that produced a composition.  A report
reads the evaluation the numbers belong to, and a count of zero means no
operator has run, which is not a flux of zero.  The evaluation matters
because the base entry is not steady while the base is: MEASURED on the
\texttt{lower\_profile} configuration, the advective half at the base face
reads $+4.47\times10^{-8}$ at evaluation 1079, $-3.35\times10^{-10}$ at
11220 and $+6.29\times10^{-10}$ at 12000, as the base sound wave turns the
face flux over in sign (item P6e).

The whole operator is gated by \texttt{He\_diffusion}, default off. With
\texttt{he\_metal\_diffusion} (default off) each trace element additionally
carries its own mass fraction and its own binary coefficient; otherwise the
metals are frozen to hydrogen, i.e.\ they ride in component 1 at the
metal/H ratio the cell holds, and what \texttt{project\_elements} preserves
for them is that ratio and not the density.

The Dirichlet composition sits at the \emph{center} of cell~1, half a cell
above the base level, which since 2026-09-27 is the lower face of the grid
(Sec.~\ref{sec:bcic:base}).  That offset is a first-order term in the cell
width which is recorded and not changed; the He/H change near the base is
$1.5\times10^{-4}$ (\texttt{md/Update\_EXHALE\_stage3.md} section 29.11).

\subsection{The interdiffusion enthalpy flux}
\label{sec:mol:interdiffusion}

Particles that move relative to the mass-averaged velocity carry their
enthalpy with them.  In a mixture whose species move with velocities
$v + w_s$, the energy equation is
%
\begin{equation}
  \frac{\partial\mathcal{E}}{\partial t}
  + \nabla\!\cdot\!\bigl[(\mathcal{E}+p)\mathbf{v}\bigr]
  = \ldots - \nabla\!\cdot q_d,
  \qquad
  q_d = \sum_s h_s J_s,\quad J_s = \rho_s w_s,\quad \sum_s J_s = 0,
  \label{eq:mol:qd}
\end{equation}
%
\citep[eqs.~11--13, the Dufour flux and the kinetic energy of the relative
motion left out there and here]{Cook2009}.  Because $J_s$ is a flux relative
to the \emph{mass}-weighted velocity, $h_s$ is the specific enthalpy and not
the internal energy.  With the binary element model of
Sec.~\ref{sec:mol:diffusion} every species moves with its element and the
electrons with their ions, so $J_1 = -J_{\rm He}$ and
%
\begin{equation}
  q_d = J_{\rm He}\,\bigl(h_{\rm He} - h_1\bigr)
      + \sum_X J_X\,\bigl(h_X - h_{\rm H}\bigr),
  \label{eq:mol:qdbinary}
\end{equation}
%
the metal sum present only where \texttt{He\_metal\_diffusion} moves each
metal against hydrogen (taken against $h_{\rm H}$, component 1 without its
metals, the group that recoils) and the metals carry mass in the mixture.
The $h_c$ are \emph{sensible} specific enthalpies built from the caloric
equation of state (\texttt{component\_specific\_enthalpies} in
\texttt{binary\_element\_diffusion.f90}): per particle
$h = e + kT = kT\,\gamma/(\gamma-1)$ for every atom, ion and molecular ion,
each ion counted with the $Z$ electrons the ambipolar field makes follow it,
and H$_2$ adding its rovibrational energy of Sec.~\ref{sec:mol:caloric};
HeH$^+$ is shared between the two components in the proportion of its
nuclei.  Formation energy is not in them: the energy update books chemical
energy through the formation-energy density, whose reference in the source
step is the composition the element transport has just written, so the
formation energy of the moved nuclei rides with them already, and the
stationary energy row holds chemical energy only as the reaction heating and
cooling inside $\mathcal{H}-\Lambda$.  $J$ is the operator's whole diffusive
flux, gradient, eddy and settling drift together, since eddy mixing carries
enthalpy as it carries the elements.  Face values of $h_c$ are the
arithmetic means of the two cells, the face rule of every coefficient of the
operator, and the flux is zero wherever $J$ is, in particular at faces 0 and
$N$.

The divergence is taken on the face areas and volumes of the element rows,
$[A_+q_d(j) - A_-q_d(j-1)]/V_j$ for $j = 2\ldots N$, so an interior face
cancels between its two cells and the energy the term moves telescopes over
the column; cell~1, the Dirichlet reservoir of the element operator, carries
no element equation and no interdiffusion energy.  Both routes carry it: the
marching step subtracts $\Delta t\,\nabla\!\cdot q_d$, formed from the element
fluxes the accepted element step itself carried, from the thermal energy
before the coupled source step starts (\texttt{EXHALE\_main.f90}); the
stationary energy row adds $\nabla\!\cdot q_d$ formed from the element fluxes
of the state itself \eqref{eq:residual}.

\paragraph{The transported carriers.}  A molecular carrier moved by the
carrier operator of Sec.~\ref{sec:mol:carrier} (H$_2$, and OH, H$_2$O, CO
with the oxygen chemistry) moves relative to the rest of its own elements as
well: its diffusive, eddy and settling number flux $\Phi_c$ moves nuclei that
the element fluxes do not, and the nuclei of the same elements held by the
species that are not carriers move the other way, nucleus for nucleus, so the
element fluxes are left as they are.  The same identity then gives
%
\begin{equation}
  q_{\rm car} = \sum_c \Phi_c\,\Bigl[h_c - \sum_{\rm el}\nu_{c,\rm el}\,
  \bar h_{\rm el}\Bigr],
  \label{eq:mol:qcar}
\end{equation}
%
with $h_c$ the enthalpy of one carrier particle and $\bar h_{\rm el}$ the
enthalpy per nucleus of el of the species that take up the recoil, each
particle with the electrons its charge gave up.  These are exactly the species
whose densities the composition update rescales (\texttt{carrier\_write\_back}),
so the energy row moves the enthalpy of the nuclei the composition rows move:
for hydrogen the species of hydrogen nuclei alone that the flow does not carry,
H\,I, H$_2^+$, H$_3^+$, and H$^+$ unless ionization transport carries it; HeH$^+$
never recoils, since rescaling it would move a helium nucleus with the
hydrogen one; for O and C their atomic stages where the metals are in the
equation of state.  An element with no recoiling species in a cell recoils
with its neutral atom.  An H$_2$ molecule rising through atomic
hydrogen at zero hydrogen flux thus carries $h_{\rm H_2} - 2h_{\rm H} =
k(u_{\rm rv} - 5T/2)$, about $-1.5\,kT$ at 1000\,K.  $\Phi_c$ is the carrier
operator's own face flux for the state in hand (\texttt{carrier\_enthalpy\_face\_flux}
in \texttt{diffusive\_photochemistry.f90}), with the temperature and particle
count of that state; both ends of the carrier column are closed, so
$q_{\rm car}(0) = q_{\rm car}(N) = 0$ and the divergence, taken over cells
$1\ldots N$ on the carrier rows' geometry, sums to zero over the column.  The
stationary energy row adds $\nabla\!\cdot q_{\rm car}$ beside $\nabla\!\cdot q_d$;
the marching step subtracts $\Delta t\,\nabla\!\cdot q_{\rm car}$ after an
accepted carrier step, from the fluxes of the composition that step returned
(its own fluxes to first order in $\Delta t$, exactly at a fixed point).  The
ionization stages carried by \texttt{Ionization transport} are not covered:
their face flux contains the element's advection, and the relative part is
not formed.  MEASURED on the certified well-mixed He/H\,$=0.19$ state of
LHS~1140\,b before the term existed: $|\nabla\!\cdot q_{\rm car}|$ is a median
1.8 per cent of the heating over $1.001 < r/R_p < 1.03$, up to 10 per cent,
comparable to the net heating where heating and cooling cancel, and that
state is no longer stationary with the term (the carrier row stands at
$7\times10^{-3}$ on the first pass).  The key
\texttt{Interdiffusion enthalpy flux:} is \textbf{on by default} whenever
\texttt{He\_diffusion} moves helium or a molecular carrier is transported, and
controls both terms, because the term belongs to the energy
equation of any mixture whose components move relative to one another;
\texttt{False} exists to reproduce published escape models that omit it
\citep[eq.~3]{Koskinen2013}, \citep[eq.~B3]{Koskinen2022}.  With
\texttt{He\_diffusion} off the element term is identically zero with no
arithmetic, and so is the carrier term without transported carriers.
MEASURED on the LHS~1140\,b \texttt{He\_diffusion} states: the divergence
reaches 13 to 26 per cent of the local energy budget at 1.1--5\,$R_p$, and on
the molecular He/H\,$=1.8$ state (24 outer passes with the term, not yet
certified) it lowers $T$ by a median 166\,K over 1.3--2\,$R_p$ and He/H by
3.1 per cent, and moves $\log_{10}\dot M$ by $+0.002$.
